\documentclass[leqno]{article}
% Shared QBT/Soln pack preamble for resources.danieljsmith.org.
\usepackage{graphicx}
\usepackage{dsfont}
\usepackage{mathtools}
\usepackage{amssymb}
\usepackage{amsmath}
\usepackage{mleftright}
\usepackage{amsthm}
\usepackage{float}
\usepackage{ragged2e}
\usepackage{tensor}
\usepackage{fancyhdr}
\usepackage{empheq}
\usepackage[most]{tcolorbox}
\usepackage{enumitem}
\usepackage{dirtytalk}
\usepackage{marginnote}
\usepackage{microtype}
\usepackage[
  a4paper,
  left=0.8in,
  right=1in,
  top=1in,
  bottom=0.5in,
  headheight=14pt,
  headsep=0.4in
]{geometry}
\setlength{\parindent}{0cm}

% TikZ / PGFPlots
\usepackage{pgfplots}
\pgfplotsset{compat=1.18}
\usepgfplotslibrary{fillbetween, polar}
\usetikzlibrary{calc, positioning}
\usetikzlibrary{angles, quotes}
\usetikzlibrary{arrows.meta}
\usetikzlibrary{decorations.pathreplacing, decorations.markings}
\usetikzlibrary{patterns}
\usetikzlibrary{intersections}

% Links and contents. The contents follow the TMUA papers: compact entries,
% no dot leaders, and each question's mark total in grey after its label.
\usepackage{tocloft}
\tocloftpagestyle{fancy}
\setlength{\cftbeforesubsecskip}{2pt}
\setlength{\cftsubsecindent}{1.5em}
\setlength{\cftsubsecnumwidth}{0pt}
\renewcommand{\cftsubsecfont}{\small\color{black}}
\renewcommand{\cftsubsecpagefont}{\small\color{black}}
\renewcommand{\cftsubsecleader}{\hfill}
\usepackage[colorlinks=true,linkcolor=black,citecolor=magenta,urlcolor=black]{hyperref}%
\urlstyle{same}

% Pack macros
% \marks also adds its value to the current question's total in the .aux
% file, so the contents show each total from the second pass. The write
% happens at shipout, so a box typeset more than once counts once. The
% pack's mark total for the contents-page title is summed the same way.
\makeatletter
\newcount\djs@totalmarks
\newcommand{\djsAddMarks}[2]{%
  \global\advance\djs@totalmarks#2\relax
  \@ifundefined{djs@marks@#1}%
    {\expandafter\xdef\csname djs@marks@#1\endcsname{#2}}%
    {\expandafter\xdef\csname djs@marks@#1\endcsname
      {\the\numexpr\csname djs@marks@#1\endcsname+#2\relax}}}
\RenewDocumentCommand{\marks}{o m}{%
  \protected@write\@auxout{}{\string\djsAddMarks{\number\value{enumi}}{#2}}%
  \IfValueTF{#1}%
    {\marginnote{\raggedright\textbf{[#2]}}[#1]}%
    {\marginnote{\raggedright\textbf{[#2]}}}%
}
\makeatother

\newcommand{\vspacerule}[2]{%
  \par\vspace*{#1}%
  \noindent\hrulefill\par
  \vspace*{#2}%
}

\definecolor{scarlet}{RGB}{205, 40, 75}
% Site colours for the contents-page title: the accent, the rule colour, a
% stronger rule colour for the path separator, and the muted text colour.
\definecolor{djsBurgundy}{RGB}{118, 59, 67}
\definecolor{djsHairline}{RGB}{216, 213, 205}
\definecolor{djsSlash}{RGB}{170, 166, 158}
\definecolor{djsMuted}{RGB}{96, 101, 106}

% Worked-solution box, adapted from the TMUA worked-solutions box. A white breakable box with a left scarlet rule and no surrounding
% frame or tint. A part that breaks fills the rest of its page, so the rule
% runs to the bottom margin; the unbroken or last part ends in a short
% horizontal foot rule so a solution visibly stops. Unlike the TMUA box there
% is no answer label. The box does not grow to the left, so the rule sits at
% the text edge; growing to the right keeps the usable line width close to the
% old tcolorbox Solution box.
% Spacing inside the box: paragraphs are 6pt apart, so source blank lines show
% (solutions put one after each display and none between a lead-in and its
% display). Every display gets 6pt above and 10pt below, whatever the length
% of its lead-in line; align rows are spread by an 8pt \jot; solution parts
% are 14pt apart and sub-parts 8pt; \left and \right space like ordinary
% brackets (mleftright); no single line is stranded at a page break; and
% neither the Solution label nor the line before a list ends a page. No
% skip can shrink, so full pages are not squeezed.
\newtcolorbox{djsSolution}{%
  enhanced, breakable, frame hidden, sharp corners, colback=white,
  height fixed for=first and middle, height fill,
  extras unbroken and last={natural height},
  borderline west={2.2pt}{0pt}{scarlet},
  left=11pt, right=0pt, top=5pt, bottom=13pt,
  before skip=13pt, after skip=4pt,
  grow to right by=10mm,
  overlay unbroken and last={%
    \fill[scarlet] (frame.south west)
      rectangle ([xshift=3.5cm, yshift=2.2pt]frame.south west);},
  before upper={{\color{scarlet}\bfseries Solution}\par\nobreak\vspace{4pt}%
    \setlength{\parskip}{6pt plus 2pt}%
    \setlength{\abovedisplayskip}{6pt plus 2pt}%
    \setlength{\abovedisplayshortskip}{6pt plus 2pt}%
    \setlength{\belowdisplayskip}{10pt plus 2pt}%
    \setlength{\belowdisplayshortskip}{10pt plus 2pt}%
    \setlength{\jot}{8pt}%
    \setlist{parsep=6pt plus 2pt, topsep=6pt plus 2pt, beginpenalty=10000}%
    \setlist[enumerate,1]{itemsep=14pt plus 2pt}%
    \setlist[enumerate,2]{itemsep=8pt plus 2pt}%
    \mleftright
    \clubpenalty=10000 \widowpenalty=10000 \displaywidowpenalty=10000\relax}}

% Optional smaller figure inside a solution box. \scalebox shrinks node text
% with the drawing. Not applied to existing figures.
\newcommand{\SolnFigureScale}{0.85}
\newsavebox{\djsSolnFigureBox}
\newenvironment{djsSolutionFigure}
  {\par\vspace{4pt}\begin{center}\begin{lrbox}{\djsSolnFigureBox}}
  {\end{lrbox}\scalebox{\SolnFigureScale}{\usebox{\djsSolnFigureBox}}%
   \end{center}\vspace{2pt}\par}

\newcommand{\cxl}[2]{%
  \tikz[baseline=(X.base)]{%
    \node[inner sep=0pt, outer sep=0pt, text=#1] (X) {$#2$};
    \draw[#1, line width=0.4pt] (X.south west)--(X.north east);
  }%
}

% Page-style hook run at the contents page and each question's opening page.
% A no-op for QBT packs; \djsSolnHeader points it at the fancy header.
\newcommand{\djsQuestionPageStyle}{}

% A contents entry: "Question N" in a box wide enough for two digits, so the
% mark totals line up, then the total in grey. Bookmarks keep "Question N".
\makeatletter
\newlength{\djsTocLabelWidth}
\newcommand{\djsTocMarks}[1]{%
  \@ifundefined{djs@marks@#1}{}{%
    {\footnotesize\itshape\color{black!45}%
     \csname djs@marks@#1\endcsname~mark%
     \ifnum\csname djs@marks@#1\endcsname=1 \else s\fi}}}
\newcommand{\djsTocEntry}[1]{%
  \texorpdfstring
    {\settowidth{\djsTocLabelWidth}{Question 00}%
     \makebox[\dimexpr\djsTocLabelWidth+0.9em\relax][l]{Question #1}%
     \djsTocMarks{#1}}%
    {Question #1}}
\makeatother

% The question count for the contents-page title is written to the .aux at
% the end of the document, so it also appears from the second pass.
\newcounter{djsQuestions}
\newcommand{\questionitem}{%
  \djsQuestionPageStyle
  \item\phantomsection
  \stepcounter{djsQuestions}%
  \addcontentsline{toc}{subsection}{\protect\djsTocEntry{\number\value{enumi}}}%
}
\makeatletter
\newcommand{\djsSetQuestionCount}[1]{\gdef\djs@questioncount{#1}}
\AtEndDocument{\immediate\write\@auxout
  {\string\djsSetQuestionCount{\number\value{djsQuestions}}}}
\makeatother

\renewcommand{\contentsname}{Questions}

% \djsPackHeader[<link content tag>]{<left header>}; the tag says which
% document a click on the header link came from.
\newcommand{\djsPackHeader}[2][questions]{%
  \pagestyle{fancy}%
  \fancyhead{}%
  \fancyfoot{}%
  \renewcommand{\headrulewidth}{0.9pt}%
  \lhead{#2}%
  \chead{}%
  \rhead{\href{https://danieljsmith.org/?utm_source=pdf&utm_campaign=further-maths&utm_content=#1}{danieljsmith.org}}%
}

% The header's topic is also the contents-page title, under which a label
% says which booklet this is.
\newcommand{\djsPackTitle}{}
\newcommand{\djsBookletLabel}{}
\newcommand{\djsQbtHeader}[1]{%
  \djsPackHeader{#1}%
  \renewcommand{\djsPackTitle}{#1}%
  \renewcommand{\djsBookletLabel}{Questions by topic}%
}

% The colophon at the foot of page one, independent of the page style: the
% contents-page title's hairline led by a short scarlet rule, then the date on
% the left and a link to the site on the right, all inside the bottom margin.
% Its sizes are fixed, so it matches the TMUA papers' colophon. The document
% supplies the date; recompiling never changes it.
\makeatletter
\newcommand{\djsLastUpdated}[1]{%
  \AtBeginDocument{\AddToHookNext{shipout/foreground}{\djs@colophon{#1}}}}
\newcommand{\djs@colophon}[1]{%
  \put(0.8in,-\dimexpr\paperheight-0.47in\relax){%
    \makebox[0pt][l]{\color{djsHairline}%
      \rule[0.9pt]{\dimexpr\paperwidth-1.8in\relax}{0.6pt}}%
    \makebox[0pt][l]{\color{scarlet}\rule{3.5cm}{2.2pt}}}%
  \put(0.8in,-\dimexpr\paperheight-0.32in\relax){\makebox[0pt][l]{%
    \fontsize{9}{11}\selectfont Last updated #1}}%
  \put(\dimexpr\paperwidth-1in\relax,-\dimexpr\paperheight-0.32in\relax){%
    \makebox[0pt][r]{\fontsize{7}{8.5}\selectfont\color{djsMuted}%
      Latest version at \href{https://resources.danieljsmith.org/}%
        {\color{black}resources.danieljsmith.org}}}}
\makeatother

% Solutions packs show the header only on the contents page and each
% question's opening page, so it never appears part-way through a solution.
% \thispagestyle marks the next page shipped out, which is the question's
% page because every solution question starts after a \newpage.
\newcommand{\djsSolnHeader}[1]{%
  \djsPackHeader[solutions]{\textit{Solutions} - #1}%
  \pagestyle{empty}%
  \renewcommand{\djsQuestionPageStyle}{\thispagestyle{fancy}}%
  \renewcommand{\djsPackTitle}{#1}%
  \renewcommand{\djsBookletLabel}{Worked solutions}%
}

% \djsSitePath{<strand>}{<topic>}, called once after the header, names where
% the pack sits on the site; the contents-page title prints it above the pack
% title, e.g. CORE PURE / FURTHER CALCULUS. Without it the line is left out.
\makeatletter
\newif\ifdjs@sitepath
\newcommand{\djs@strand}{}
\newcommand{\djs@section}{}
\newcommand{\djsSitePath}[2]{%
  \djs@sitepathtrue
  \renewcommand{\djs@strand}{#1}%
  \renewcommand{\djs@section}{#2}%
}
\makeatother

% Further Mechanics packs call \djsMechanicsPack in their preamble, so that
% \djsFrontMatter prints the default value of g above the contents. The box
% keeps its own space above it (before skip is the whole gap below the title,
% not an addition to it). Without the box, \djsGravityNote is just extra space
% before the contents heading, bringing the gap above it close to the gap
% between the box and the heading on Further Mechanics pages.
\newcommand{\djsGravityNote}{\vspace{14pt}}
\newcommand{\djsMechanicsPack}{%
  \renewcommand{\djsGravityNote}{%
    \begin{tcolorbox}[sharp corners, colback=white, colframe=black,
      boxrule=1pt, left=8pt, right=8pt, top=6pt, bottom=6pt,
      before skip=20pt]
    {\bfseries Instructions}\par\vspace{4pt}
    Unless a question states otherwise, take $g = 9.8\,\mathrm{m\,s^{-2}}$
    wherever a numerical value of $g$ is required. Solutions also give final
    answers for $g = 10\,\mathrm{m\,s^{-2}}$ where these differ.
    \end{tcolorbox}%
    \vspace{16pt}}%
}

% Contents-page title: the site path in spaced burgundy capitals, the pack
% title, a hairline led by a short scarlet rule (the solution box's foot rule),
% then the booklet label with the question and mark counts. The counts come
% from the .aux, so they appear from the second pass.
\makeatletter
\newsavebox{\djs@eyebrowbox}
\newcommand{\djs@caps}[1]{\textls[180]{\MakeUppercase{#1}}}
\newcommand{\djs@eyebrow}{%
  \color{djsBurgundy}\footnotesize\bfseries
  \djs@caps{\djs@strand}%
  {\color{djsSlash}\hspace{0.55em}/\hspace{0.55em}}%
  \djs@caps{\djs@section}}
\newcommand{\djs@dot}{\hspace{0.45em}\textperiodcentered\hspace{0.45em}}
\newcommand{\djs@counts}{%
  \@ifundefined{djs@questioncount}{}{%
    \djs@questioncount~question\ifnum\djs@questioncount=1 \else s\fi
    \ifnum\djs@totalmarks>0
      \djs@dot\the\djs@totalmarks~mark\ifnum\djs@totalmarks=1 \else s\fi
    \fi}}
\newcommand{\djsContentsTitle}{%
  \ifdjs@sitepath
    \sbox\djs@eyebrowbox{\djs@eyebrow}%
    \ifdim\wd\djs@eyebrowbox>\linewidth
      \@latex@warning@no@line{The \string\djsSitePath\space line is wider
        than the text, so it wraps. Shorten its topic}%
    \fi
    {\raggedright\djs@eyebrow\par}%
    \vspace{7pt}%
  \fi
  {\raggedright\fontsize{24.88}{28}\selectfont\djsPackTitle\par}%
  \vspace{9pt}%
  \noindent\makebox[0pt][l]{\color{djsHairline}\rule[0.9pt]{\linewidth}{0.6pt}}%
  {\color{scarlet}\rule{3.5cm}{2.2pt}}\par
  \vspace{5pt}%
  {\small\color{djsMuted}\djsBookletLabel\hfill\djs@counts\par}%
  \vspace{8pt}%
}
\makeatother

\newcommand{\djsFrontMatter}{%
  \djsQuestionPageStyle
  \vspace*{6pt}%
  \djsContentsTitle
  \djsGravityNote
  \tableofcontents
  \newpage
}


\djsQbtHeader{Eigenvalues, Eigenvectors and Diagonalisation}
\djsSitePath{Core Pure}{Vectors, Matrices \& Linear Transformations}
\djsLastUpdated{27 October 2026}

\begin{document}
\djsFrontMatter
\begin{enumerate}[
  label=\textbf{\arabic*.},
  leftmargin=*,
  topsep=0pt,
  partopsep=0pt,
  parsep=0pt
]

% ---- Question 1 ----
\questionitem
% [Source: _QBT__Eigenvalues__Eigenvectors_and_Diagonalisation, Q1]

The matrix $\mathbf{A}$ is given by
\[\mathbf{A} = \begin{pmatrix} 3 & k \\ 2 & -2 \end{pmatrix}\]
where $k$ is a constant. \\

Given that $4$ is an eigenvalue of $\mathbf{A}$, find
\begin{itemize}
  \item the value of $k$,
  \item the other eigenvalue of $\mathbf{A}$,
  \item an eigenvector associated with each eigenvalue of $\mathbf{A}$.
\end{itemize}
\marks[-10pt]{6}

\newpage

% ---- Question 2 ----
\questionitem
% [Source: _QBT__Eigenvalues__Eigenvectors_and_Diagonalisation, Q2]

The matrix $\mathbf{M}$ is
\[\mathbf{M} = \begin{pmatrix} a & 0.6 \\ 0.6 & b \end{pmatrix}\]
where $a$ and $b$ are constants with $a > 0$. \\

The matrix $\mathbf{M}$ represents a reflection in a line through the origin. \\
\begin{enumerate}[label=\textbf{(\alph*)}]
  \item Write down the eigenvalues of $\mathbf{M}$. \marks{1} \\
  \item By considering the characteristic equation of $\mathbf{M}$, use your answer to part (a) to find the values of $a$ and $b$. \marks{3} \\
  \item Hence find the equation of the mirror line of the reflection represented by $\mathbf{M}$. \marks{2} \\
\end{enumerate}

\newpage

% ---- Question 3 ----
\questionitem
% [Source: _QBT__Eigenvalues__Eigenvectors_and_Diagonalisation, Q3]

The matrix $\mathbf{M}$ is given by
\[\mathbf{M} = \begin{pmatrix} 2 & k \\ k & 5 \end{pmatrix}\]
where $k$ is a negative constant. \\

It is given that $1$ is an eigenvalue of $\mathbf{M}$. \\

\begin{enumerate}[label=\textbf{(\alph*)}]
  \item Find the value of $k$. Hence find the other eigenvalue of $\mathbf{M}$, and a corresponding eigenvector for each eigenvalue of $\mathbf{M}$. \marks{6} \\
  \item
  \begin{enumerate}[label=(\roman*)]
    \item Describe how the directions of the invariant lines through the origin of the transformation represented by $\mathbf{M}$ are related to each other. Fully justify your answer. \marks{2}
    \item Describe fully the transformation represented by $\mathbf{M}$, stating the equation of any line of invariant points. \marks{3} \\
  \end{enumerate}
\end{enumerate}

\newpage

% ---- Question 4 ----
\questionitem
% [Source: _QBT__Eigenvalues__Eigenvectors_and_Diagonalisation, Q4]

A $3 \times 3$ matrix $\mathbf{A}$ has eigenvalues $2$, $-1$ and $1$, with corresponding eigenvectors
\[\begin{pmatrix} 1 \\ 0 \\ 1 \end{pmatrix}, \quad \begin{pmatrix} 0 \\ 1 \\ 1 \end{pmatrix}, \quad \begin{pmatrix} 1 \\ 1 \\ 1 \end{pmatrix}\]
respectively. \\
\begin{enumerate}[label=\textbf{(\alph*)}]
  \item By expressing $\mathbf{A}$ in the form $\mathbf{P}\mathbf{D}\mathbf{P}^{-1}$, where $\mathbf{D}$ is a diagonal matrix, find the matrix $\mathbf{A}$. \marks{4} \\
  \item Hence find the matrix $\mathbf{A}^4$. \marks{2} \\
\end{enumerate}

\newpage

% ---- Question 5 ----
\questionitem
% [Source: _QBT__Eigenvalues__Eigenvectors_and_Diagonalisation, Q5]

The matrix $\mathbf{M}$ is given by
\[\mathbf{M} = \begin{pmatrix} 4 & -1 & 1 \\ a & 5 & a \\ -8 & b & -5 \end{pmatrix}\]
where $a$ and $b$ are real numbers. \\

The vectors
\[\mathbf{v}_1 = \begin{pmatrix} 1 \\ 2 \\ 0 \end{pmatrix} \quad \text{and} \quad \mathbf{v}_2 = \begin{pmatrix} 0 \\ 1 \\ 1 \end{pmatrix}\]
are eigenvectors of $\mathbf{M}$, with corresponding eigenvalues $\lambda_1$ and $\lambda_2$ respectively. \\

\begin{enumerate}[label=\textbf{(\alph*)}]
  \item
  \begin{enumerate}[label=(\roman*)]
    \item Find the value of $\lambda_1$. \marks{1}
    \item Find the value of $a$ and the value of $b$. \marks{2} \\
  \end{enumerate}
\end{enumerate}
It is also given that $\lambda_3 = 3$ is an eigenvalue of $\mathbf{M}$. \\
\begin{enumerate}[label=\textbf{(\alph*)}, resume]
  \item Find the value of $\lambda_2$, and find an eigenvector $\mathbf{v}_3$ of $\mathbf{M}$ corresponding to the eigenvalue $\lambda_3 = 3$. \marks{4} \\
  \item Find matrices $\mathbf{U}$, $\mathbf{D}$ and $\mathbf{U}^{-1}$, such that $\mathbf{D}$ is diagonal and $\mathbf{M} = \mathbf{U}\mathbf{D}\mathbf{U}^{-1}$. \marks{3} \\
\end{enumerate}

\newpage

% ---- Question 6 ----
\questionitem
% [Source: _QBT__Eigenvalues__Eigenvectors_and_Diagonalisation, Q6]

\[\mathbf{M} = \begin{pmatrix} 1 & -2 & 2 \\ -2 & 2 & 0 \\ 2 & 0 & 0 \end{pmatrix}\] \\
\begin{enumerate}[label=\textbf{(\alph*)}]
  \item Show that $-2$ is an eigenvalue of $\mathbf{M}$, and find a corresponding eigenvector. \marks{4} \\
  \item Given that $\mathbf{MU} = \mathbf{UD}$, where $\mathbf{D}$ is a diagonal matrix and each column of $\mathbf{U}$ is an eigenvector of $\mathbf{M}$ of unit length, find possible matrices for $\mathbf{D}$ and $\mathbf{U}$. \marks{11} \\
\end{enumerate}

\newpage

% ---- Question 7 ----
\questionitem
% [Source: _QBT__Eigenvalues__Eigenvectors_and_Diagonalisation, Q7]

The matrix $\mathbf{A}$ is given by
\[\mathbf{A} = \begin{pmatrix} \dfrac{3}{8} & \dfrac{5}{16} \\[6pt] \dfrac{5}{4} & \dfrac{3}{8} \end{pmatrix}\] \\
\begin{enumerate}[label=\textbf{(\alph*)}]
  \item Find the eigenvalues of $\mathbf{A}$ and a corresponding eigenvector for each eigenvalue. \marks{5} \\
  \item Write down a matrix $\mathbf{P}$ and a diagonal matrix $\mathbf{D}$ such that $\mathbf{A} = \mathbf{P}\mathbf{D}\mathbf{P}^{-1}$. \marks{1} \\
  \item Find $\mathbf{A}^n$, where $n$ is a positive integer, as a single matrix whose entries are expressed in terms of $n$. \\
  Hence find the matrix $\mathbf{L}$ such that $\mathbf{A}^n \to \mathbf{L}$ as $n \to \infty$, explaining briefly why this limit exists. \marks{5} \\
  \item The transformation represented by $\mathbf{L}$ maps every point of the plane onto a single line through the origin. \\
  Find the equation of this line. \marks{2} \\
\end{enumerate}

\newpage

% ---- Question 8 ----
\questionitem
% [Source: _QBT__Eigenvalues__Eigenvectors_and_Diagonalisation, Q8]

The matrix $\mathbf{A}$ is given by
\[\mathbf{A} = \begin{pmatrix} -4 & -13 & 8 \\ -6 & -7 & 6 \\ -12 & -22 & 16 \end{pmatrix}\]
\begin{enumerate}[label=\textbf{(\alph*)}]
  \item Show that the characteristic equation of $\mathbf{A}$ is $\lambda^3 - 5\lambda^2 + 2\lambda + 8 = 0$ and hence find the eigenvalues of $\mathbf{A}$. \marks{4} \\
  \item Find a matrix $\mathbf{P}$ and a diagonal matrix $\mathbf{D}$ such that $\mathbf{A} = \mathbf{P}\mathbf{D}\mathbf{P}^{-1}$. \marks{5} \\
  \item Hence show that, for all positive integers $n$, the element in the first row and first column of $\mathbf{A}^n$ is $2(-1)^n - 2^n$. \marks{3} \\
\end{enumerate}

\newpage

% ---- Question 9 ----
\questionitem
% [Source: _QBT__Eigenvalues__Eigenvectors_and_Diagonalisation, Q9]

The matrix $\mathbf{M}$ is given by
\[\mathbf{M} = \begin{pmatrix} 3 & 4 & -4 \\ 1 & 3 & -1 \\ 1 & 4 & -2 \end{pmatrix}\]
\begin{enumerate}[label=\textbf{(\alph*)}]
  \item Show that the characteristic equation of $\mathbf{M}$ is
  \[\lambda^3 - 4\lambda^2 + \lambda + 6 = 0\] \marks[-20pt]{3}
  \item Hence find the three eigenvalues of $\mathbf{M}$. \marks{2} \\
  \item For each eigenvalue of $\mathbf{M}$ find an associated eigenvector. \marks{4} \\
  \item Write down a matrix $\mathbf{P}$ and a diagonal matrix $\mathbf{D}$ such that $\mathbf{M} = \mathbf{P}\mathbf{D}\mathbf{P}^{-1}$. Hence find $\mathbf{M}^n$ as a single matrix, where $n$ is a positive integer, giving each entry in terms of $n$. \marks{4} \\
\end{enumerate}

\newpage

% ---- Question 10 ----
\questionitem
% [Source: _QBT__Eigenvalues__Eigenvectors_and_Diagonalisation, Q10]

The square matrix $\mathbf{A}$ has an eigenvalue $\lambda$ with corresponding eigenvector $\mathbf{e}$. The non-singular matrix $\mathbf{M}$ is of the same order as $\mathbf{A}$, and the matrix $\mathbf{B}$ is defined by $\mathbf{B} = \mathbf{M}^{-1}\mathbf{A}\mathbf{M}$. \\
\begin{enumerate}[label=\textbf{(\alph*)}]
  \item Show that $\mathbf{M}^{-1}\mathbf{e}$ is an eigenvector of $\mathbf{B}$, and that $\lambda$ is the corresponding eigenvalue. \marks{3} \\
\end{enumerate}
Let
\[\mathbf{A} = \begin{pmatrix} 3 & 0 & 0 \\ 1 & -2 & 0 \\ 2 & 5 & 1 \end{pmatrix}\]
\begin{enumerate}[label=\textbf{(\alph*)}, resume]
  \item Write down the eigenvalues of $\mathbf{A}$ and obtain corresponding eigenvectors. \marks{4} \\
\end{enumerate}
It is given that
\[\mathbf{M} = \begin{pmatrix} 1 & 1 & 0 \\ 0 & 1 & 0 \\ 0 & 2 & 1 \end{pmatrix}\]
\begin{enumerate}[label=\textbf{(\alph*)}, resume]
  \item Find $\mathbf{M}^{-1}$ and hence, without finding $\mathbf{B}$ explicitly, find the eigenvalues and corresponding eigenvectors of $\mathbf{B}$. \marks{5} \\
\end{enumerate}

\newpage

% ---- Question 11 ----
\questionitem
% [Source: _QBT__Eigenvalues__Eigenvectors_and_Diagonalisation, Q11]

It is given that $\lambda$ is an eigenvalue of the non-singular square matrix $\mathbf{A}$, with corresponding eigenvector $\mathbf{e}$. \\
\begin{enumerate}[label=\textbf{(\alph*)}]
  \item Show that $\lambda^{2}$ is an eigenvalue of $\mathbf{A}^{2}$ and that $\lambda^{-1}$ is an eigenvalue of $\mathbf{A}^{-1}$, and that $\mathbf{e}$ is a corresponding eigenvector in each case. \marks{3} \\
  \item Deduce that $\lambda^{2} + 2\lambda^{-1}$ is an eigenvalue of $\mathbf{A}^{2} + 2\mathbf{A}^{-1}$. \marks{1} \\
\end{enumerate}
It is given that $-1$ is an eigenvalue of the matrix $\mathbf{A}$, where
\[\mathbf{A} = \begin{pmatrix} 2 & 1 & -1 \\ 3 & 2 & -3 \\ 3 & 1 & -2 \end{pmatrix}\]
\begin{enumerate}[label=\textbf{(\alph*)}, resume]
  \item Find a corresponding eigenvector. \marks{2} \\
  \item It is also given that
  \[\begin{pmatrix} 1 \\ 0 \\ 1 \end{pmatrix} \quad \text{and} \quad \begin{pmatrix} 1 \\ 1 \\ 1 \end{pmatrix}\]
  are eigenvectors of the matrix $\mathbf{A}$. Find the corresponding eigenvalues. \marks{2} \\
  \item Hence find a matrix $\mathbf{P}$ and a diagonal matrix $\mathbf{D}$ such that
  \[\left(\mathbf{A}^{2} + 2\mathbf{A}^{-1}\right)^2 = \mathbf{P}\mathbf{D}\mathbf{P}^{-1}\] \marks[-20pt]{4}
\end{enumerate}

\newpage

% ---- Question 12 ----
\questionitem
% [Source: _QBT__Eigenvalues__Eigenvectors_and_Diagonalisation, Q12]

The matrix $\mathbf{M}$ is given by
\[\mathbf{M} = \dfrac{1}{3}\begin{pmatrix} 2 & -1 & 2 \\ 2 & 2 & -1 \\ -1 & 2 & 2 \end{pmatrix}\] \\
\begin{enumerate}[label=\textbf{(\alph*)}]
  \item
  \begin{enumerate}[label=(\roman*)]
    \item Find the characteristic equation of $\mathbf{M}$, giving your answer in the form $\lambda^3 + p\lambda^2 + q\lambda + r = 0$, where $p$, $q$ and $r$ are integers. Hence show that $1$ is an eigenvalue of $\mathbf{M}$. \marks{3}
    \item Find the other two eigenvalues of $\mathbf{M}$, giving each in the form $x + yi$ where $x$ and $y$ are real. Deduce that $1$ is the only real eigenvalue of $\mathbf{M}$. \marks{3} \\
  \end{enumerate}
  \item A student claims that $\mathbf{M}$ can be written in the form $\mathbf{M} = \mathbf{P}\mathbf{D}\mathbf{P}^{-1}$, where $\mathbf{D}$ is a diagonal matrix and \textbf{both} $\mathbf{P}$ and $\mathbf{D}$ have only real entries. Using your answer to part (a)(ii), explain why the student is wrong. \marks{2} \\
\end{enumerate}
You are now given the following information.
\begin{itemize}
  \item $\mathbf{M}$ represents a rotation in 3-D space by an acute angle about an axis through the origin.
  \item $\mathbf{e}$ is a real eigenvector of $\mathbf{M}$ associated with the eigenvalue $1$.
\end{itemize}
\begin{enumerate}[label=\textbf{(\alph*)}, resume]
  \item
  \begin{enumerate}[label=(\roman*)]
    \item Find a possible vector $\mathbf{e}$. \marks{2}
    \item Explain the significance of $\mathbf{e}$, and the fact that its associated eigenvalue is $1$, in relation to the rotation that $\mathbf{M}$ represents. Hence write down a vector equation for the axis of this rotation. \marks{2}
    \item Explain \textbf{geometrically} why $\mathbf{M}$ has no real eigenvectors other than non-zero multiples of $\mathbf{e}$. \marks{1} \\
  \end{enumerate}
\end{enumerate}

\newpage

% ---- Question 13 ----
\questionitem
% [Source: _QBT__Eigenvalues__Eigenvectors_and_Diagonalisation, Q13]

The matrix $\mathbf{A}$ is given by
\[\mathbf{A} = \begin{pmatrix} 6 & -2 & 1 \\ -2 & 3 & 2 \\ 1 & 2 & 6 \end{pmatrix}\]
\begin{enumerate}[label=\textbf{(\alph*)}]
  \item Find the characteristic equation of $\mathbf{A}$. Hence show that $1$ is an eigenvalue of $\mathbf{A}$, and that $\mathbf{A}$ has one other eigenvalue, which is repeated. State the value of this repeated eigenvalue. \marks{5} \\
  \item Find an eigenvector of $\mathbf{A}$ corresponding to the eigenvalue $1$. Find also two eigenvectors of $\mathbf{A}$, corresponding to the repeated eigenvalue, which are perpendicular to each other. \marks{6} \\
  \item Hence find an orthogonal matrix $\mathbf{P}$ and a diagonal matrix $\mathbf{D}$ such that $\mathbf{P}^{\mathrm{T}}\mathbf{A}\mathbf{P} = \mathbf{D}$. \marks{3} \\
\end{enumerate}

\newpage

% ---- Question 14 ----
\questionitem
% [Source: _QBT__Eigenvalues__Eigenvectors_and_Diagonalisation, Q14]

\[\mathbf{A} = \begin{pmatrix} -3 & 2 & -2 \\ -2 & 2 & p \\ 10 & -5 & 3 \end{pmatrix} \quad \text{where } p \text{ is a constant}\] \\
Given that $\begin{pmatrix} 1 \\ 2 \\ -1 \end{pmatrix}$ is an eigenvector for $\mathbf{A}$, \\
\begin{enumerate}[label=\textbf{(\alph*)}]
  \item
  \begin{enumerate}[label=(\roman*)]
    \item determine the eigenvalue corresponding to this eigenvector, \marks{1}
    \item hence show that $p = -4$, \marks{2}
    \item determine the remaining eigenvalues and corresponding eigenvectors of $\mathbf{A}$. \marks{7} \\
  \end{enumerate}
  \item Write down a matrix $\mathbf{P}$ and a diagonal matrix $\mathbf{D}$ such that $\mathbf{P}^{-1}\mathbf{A}\mathbf{P} = \mathbf{D}$. \marks{1} \\
  \item
  \begin{enumerate}[label=(\roman*)]
    \item Solve the differential equation $\dot{u} = ku$, where $k$ is a constant. \marks{2} \\
    With respect to a fixed origin $O$, the velocity of a particle moving through space is modelled by
    \[\begin{pmatrix} \dot{x} \\ \dot{y} \\ \dot{z} \end{pmatrix} = \mathbf{A} \begin{pmatrix} x \\ y \\ z \end{pmatrix}\]
    By considering
    \[\begin{pmatrix} u \\ v \\ w \end{pmatrix} = \mathbf{P}^{-1} \begin{pmatrix} x \\ y \\ z \end{pmatrix} \quad \text{so that} \quad \begin{pmatrix} \dot{u} \\ \dot{v} \\ \dot{w} \end{pmatrix} = \mathbf{P}^{-1} \begin{pmatrix} \dot{x} \\ \dot{y} \\ \dot{z} \end{pmatrix}\]
    \item determine a general solution for the displacement of the particle. \marks{4} \\
    At time $t = 0$ the particle is at the point with position vector $\begin{pmatrix} 3 \\ 7 \\ -1 \end{pmatrix}$. \\
    \item Find expressions for $x$, $y$ and $z$ in terms of $t$ for this particle. \marks{2} \\
  \end{enumerate}
\end{enumerate}

\end{enumerate}
\end{document}
